\documentclass[a4paper]{article}
\usepackage{amsmath}
\usepackage{amsfonts}
\usepackage{amssymb}
\usepackage{graphicx}
\usepackage{float}

\begin{document}

\noindent \large Auteur: Siebe Haché \\
\noindent \large Studierichting: Informatica \\
\noindent \large Oefeningenreeks 6D nr. 31.29\\

\medskip
\normalsize

Onderzoek de convergentie van de reeks $\displaystyle\sum_{n=0}^{\infty}\dfrac{(-1)^n7^n}{n!}$ \\

d'Alembert\\

Stel $x_n = \dfrac{(-1)^n 7^n}{n!}$, dan is $|x_n| = \dfrac{7^n}{n!}$\\

Stel $|x_{n+1}| = \dfrac{7^{n+1}}{(n+1)!}$\\

$L = \lim\limits_{n \to \infty} \left| \dfrac{x_{n+1}}{x_n} \right|$\\

$= \lim\limits_{n \to \infty} \left( \dfrac{7^{n+1}}{(n+1)!} \cdot \dfrac{n!}{7^n} \right)$ \\

$= \lim\limits_{n \to \infty} \left( \dfrac{7 \cdot 7^n}{(n+1) \cdot n!} \cdot \dfrac{n!}{7^n} \right) $\\

$= \lim\limits_{n \to \infty} \dfrac{7}{n+1}$\\

$= 0$\\

Aangezien $L = 0 < 1$, is de reeks absoluut convergent.\\

\begin{thebibliography}{999}
\end{thebibliography}

\end{document}